% Part: model-theory
% Chapter: models-of-arithmetic
% Section: introduction

\documentclass[../../../include/open-logic-section]{subfiles}

\begin{document}

\olfileid{mod}{mar}{int}
\section{مقدمه}

\emph{مدل استاندارد} حساب همان !!{structure}~$\Struct{N}$ با
$\Domain{N} = \Nat$ است که در آن $\Obj{0}$، $\prime$، $+$، $\times$ و
$<$ آن‌گونه که انتظار می‌رود تعبیر می‌شوند. یعنی $\Obj{0}$ همان~$0$
است، $\prime$ تابع جانشین است، $+$ به جمع و $\times$ به ضرب اعداد
در~$\Nat$ تعبیر می‌شود. به‌طور مشخص،
\begin{align*}
  \Assign{\Obj{0}}{N} & = 0\\
  \Assign{\prime}{N}(n) & = n + 1\\
  \Assign{+}{N}(n, m) & = n + m\\
  \Assign{\times}{N}(n, m) & = nm\\
  \Assign{<}{N} & =
    \Setabs{\tuple{n,m}}{n,m\in\Nat \text{ و } n<m}
\end{align*}
البته برای~$\Lang{L_A}$ ساختارهایی با دامنه‌هایی جز~$\Nat$ نیز
وجود دارند. برای نمونه، می‌توانیم $\Struct{M}$ را با دامنهٔ
$\Domain{M} = \{a\}^*$ (دنباله‌های متناهیِ تک‌نماد~$a$، یعنی
$\emptyset$، $a$، $aa$، $aaa$، \dots) و تعبیرهای زیر در نظر بگیریم:
\begin{align*}
  \Assign{\Obj{0}}{M} & = \emptyset\\
  \Assign{\prime}{M}(s) & = s \concat a\\
  \Assign{+}{M}(a^n, a^m) & = a^{n + m}\\
  \Assign{\times}{M}(a^n, a^m) & = a^{nm}\\
  \Assign{<}{M} & =
    \Setabs{\tuple{a^n,a^m}}{n,m\in\Nat \text{ و } n<m}
\end{align*}
این دو ساختار از این حیث «اساساً یکسان»اند که تنها !!{element}s
!!{domain}s آن‌ها متفاوت است، نه چگونگی ارتباط !!{element}s
!!{domain}s آن‌ها با یکدیگر به‌وسیلهٔ توابع تعبیر. می‌گوییم این دو
!!{structure}s
\emph{همریخت}اند.

یکی از پیامدهای سادهٔ قضیهٔ فشردگی این است که هر نظریه‌ای که
در~$\Struct{N}$ صادق باشد، مدل‌هایی نیز دارد که با~$\Struct{N}$
همریخت نیستند. چنین ساختارهایی را
\emph{نااستاندارد} می‌نامیم. نکتهٔ جالب دربارهٔ آن‌ها این است که
!!{element}s یک مدل استاندارد (یعنی $\Struct{N}$ و نیز هر
!!{structure}s همریخت با آن) همگی با مقادیر عددنماهای استاندارد
$\num{n}$ پوشش داده می‌شوند؛ یعنی
\[
\Domain{N} = \Setabs{\Value{\num{n}}{N}}{n \in \Nat}.
\]
اما در مدل‌های نااستاندارد چنین نیست: اگر $\Struct{M}$ نااستاندارد
باشد، دست‌کم یک $x \in \Domain{M}$ وجود دارد چنان‌که برای همهٔ~$n$،
$x \neq \Value{\num{n}}{M}$.

این عناصر نااستاندارد بسیار جالب‌اند: آن‌ها «اعداد طبیعی نامتناهی»اند.
اما وجودشان، از یک جهت، پدیده‌های ناتمامیت را نیز توضیح می‌دهد. برای
نمونه گزارهٔ سازگاری حساب پئانو، $\OCon[\Th{PA}]$، یعنی
$\lnot \lexists[x][\OPrf[\Th{PA}](x, \gn{\lfalse})]$ را در نظر
بگیرید. چون $\Th{PA}$ نه $\OCon[\Th{PA}]$ را ثابت می‌کند و نه
$\lnot \OCon[\Th{PA}]$ را، هر یک را می‌توان به‌طور سازگار به
$\Th{PA}$ افزود. چون $\Th{PA}$ سازگار است،
$\Sat{N}{\OCon[\Th{PA}]}$ و در نتیجه
$\Sat/{N}{\lnot \OCon[\Th{PA}]}$. پس $\Struct{N}$
\emph{مدلِ} $\Th{PA} \cup \{\lnot \OCon[\Th{PA}]\}$ نیست و همهٔ
مدل‌های این نظریه باید نااستاندارد باشند. مدل‌های
$\Th{PA} \cup \{\lnot \OCon[\Th{PA}]\}$ باید !!{element}ی داشته
باشند که به‌منزلهٔ شاهد،
$\lexists[x][\OPrf[\Th{PA}](x, \gn{\lfalse})]$ را صادق کند؛ یعنی
عدد گودلِ !!a{derivation} از تناقض بر پایهٔ~$\Th{PA}$. چنین
!!a{element}ی نمی‌تواند استاندارد باشد، زیرا برای هر~$n$،
$\Th{PA} \Proves \lnot \OPrf[\Th{PA}](\num{n}, \gn{\lfalse})$.

\end{document}
